\section{基本范畴}
\begin{enumerate}
  \item 6 个核心概念。每年至少考其中一个概念。
  \item 目的。
        \begin{enumerate}
          \item 根本目的 \yjt 维护社会秩序。直接目的 \yjt 惩戒犯罪，保障人权。
          \item 理论分类。
                \begin{enumerate}
                  \item 犯罪控制模式和正当程序模式。
                  \item 家庭模式。
                  \item 实体真实主义。强调发现真相。宁可错放也不冤判 \yjt 消极|不放过任何一个 \yjt 积极。
                  \item 正当程序主义。强调过程。
                \end{enumerate}
        \end{enumerate}
  \item 价值。每一种价值都有工具、独立 2 个层面。侧重刑事诉讼对国家、社会、公民的效用和意义。
        \begin{enumerate}
          \item 公正（核心地位）。程序 + 实体。
          \item 秩序。社会秩序 + 追究犯罪活动有序 + 依程序法办事。
          \item 效益。效率 + 效益。
          \item 三者相互依存、相互作用、相互制约、不可偏废。
          \item \al{效益价值属于刑诉的工具价值，而不属刑诉的独立价值 \cuo}。
        \end{enumerate}
  \item 主体。扮演啥角色。
        \begin{enumerate}
          \item 专门机关。3 大 4 小。公（+ 国安、监狱、军队、海警局）检法 3 大机关。
          \item 诉讼参与人。2 种人。
                \begin{enumerate}
                  \item 当事人。
                        \begin{enumerate}
                          \item 公诉。被害人，嫌疑人/被告人。
                          \item 自诉。自诉人，被告人。
                          \item 附民。原告人，被告人。
                        \end{enumerate}
                  \item 其他诉讼参与人。法诉辩证翻鉴。法定代理人、诉讼代理人、辩护人、证人、翻译、鉴定人。
                \end{enumerate}
        \end{enumerate}
  \item 职能。各个主体在刑诉中的职责、作用和功能。控辩审。严格区分主体。
        \begin{enumerate}
          \item 控诉。公诉人、自诉人、被害人及其法定代理人、诉讼代理人。仅针对刑事而言。
                这个行为成立，导致不能追究他的刑事责任，那么这个行为就不属于控诉行为。
                \al{排非 \cuo。提出有利的量刑证据 \dui}。
          \item 辩护。嫌疑人/被告人、辩护人、法定代理人。维护被告合法权益。
          \item 审判。法院。居中裁判。
          \item 不是所有刑诉主体都有刑诉职能。\al{证人、翻译、鉴定人}。
          \item 带控方前缀、坐在控诉席的未必是控诉职能。\al{控方证人、附带民事原告人}。
          \item 仅在审判阶段，三方才共同存在。
        \end{enumerate}
  \item 构造。控辩审三方在刑诉中的格局、地位、关系。
        \begin{enumerate}
          \item 目的决定构造、价值观（潜移默化）影响构造。
          \item 审判要中立，控审要分离，控辩要平等对抗。等腰三角形。
                \begin{enumerate}
                  \item 控审要分离。没有控诉主体提起，法院就不能主动。
                  \item \al{建议、意见 \yjt 法律监督权}。
                \end{enumerate}
          \item 类型。
                \begin{enumerate}
                  \item 当事人主义。当事人居主导地位。英美法系。选项表现为弱化法官。\al{认罪认罚}。
                  \item 职权主义（我国）。出现法官。侧重实体真实。大陆法系。
                  \item 混合式。职权主义吸收当事人主义，然后反转为以当事人为主。反转为人。\al{日本、意大利}。
                \end{enumerate}
        \end{enumerate}
  \item 阶段。程序有先后，可跳跃、回转、天折。
        \begin{enumerate}
          \item 划分标准。直接任务、参加诉讼的机关和个人的构成、诉讼行为方式、诉讼法律关系、诉讼文书。
          \item 公诉。立案、侦查、起诉、审判、执行。
          \item 监察委。立案、调查、起诉、审判、执行。依监察法，严格来讲不属于刑事诉讼。
          \item 自诉。起诉、审判、执行。
        \end{enumerate}
\end{enumerate}
